(1) 水平矩形明渠内，水密度 \(\rho\)、重力加速度 \(g\)，从上游水深 \(h_1\)、平均流速 \(u_1>0\) 进入定常水跃，下游为 \(h_2,u_2\)。忽略局部床面摩擦的动量外力并采用静水压力剖面。由质量流率和动量通量推出共轭水深方程。

(2) 设上游Froude数 \(Fr_1=u_1/\sqrt{gh_1}>1\)，求物理解 \(h_2/h_1\)，证明水跃使水深增加而流速减小。

(3) 求上下游比能水头 \(H_j=h_j+u_j^2/(2g)\) 的差，证明 \(H_1-H_2=(h_2-h_1)^3/(4h_1h_2)>0\)。

(4) 从水跃下游Froude数验证流态由超临界转为亚临界，讨论 \(Fr_1\downarrow1\) 的弱水跃及能量耗散极限。
